\section{A fixed degree filter is compatible with both phases}
\label{sec:revisit-pruning}
This is an optional extension. The fixed-degree observation already appears
in the research notes of \texttt{deterministic\_op2\_20260905}; the energy
and moving-comparator checks below adapt it to this new recurrence.

Fix $\rho=\epsppr/2$ and define analytically
$A=\{i:d_i\leq1/\rho\}$. Do not enumerate $A$. Classify a newly exposed
label using its degree reply, and permanently omit it from kinetic reports
when it lies outside $A$. In particular, its adjacency list is never opened.
All formulas and row charges retain original degrees. This is the
principal-matrix algorithm with $\mM=\mQ_{AA}$, implemented lazily.

\begin{proposition}[Fixed pruning preserves the guarantees]
\label{prop:revisit-fixed-pruning}
Both the exact-real and the directed-rounded two-phase algorithms retain
their semantic guarantees and their respective constants $44$ and $60$
under this fixed filter. Every row they open has original degree at most
$1/\rho$.
\end{proposition}
\begin{proof}
For every actual $r\geq\rho$, the full RPPR optimum is supported in $A$
by $\vol(\supp\vx_r^*)\leq1/r$. Thus the restricted optimum is the
full optimum extended by zero. The nonzero branch's seed belongs to $A$.
Write $\vomega=\vomega_A$, $\vb=\vb_A$, and
$\vt=\mM^{-1}\vb$ within this proof. Principal-matrix properties give
\[
 \alpha\mI\preceq\mM\preceq\mI,\qquad \mM^{-1}\geq0,
 \qquad \mM\vomega\geq\alpha\vomega,
 \qquad |\mM|\vomega\leq\vomega.
\]
The required auxiliary objective is now
\[
 \mathcal A_\lambda(\vx)=\tfrac12\|\mM\vx-\vb\|_2^2
                            +\lambda(\mM\vomega)^\trans\vx.
\]
Its $\mM$-metric gradient is $\mM\vx-\vb+\lambda\vomega$,
and $\mathcal A_\lambda(\vt)=\lambda\vomega^\trans\vb=\alpha\lambda$.
The orthant sector uses only the Stieltjes signs and $\vb\geq0$.
Changing $\lambda$ changes the bank by
\[
 (\lambda^+-\lambda)
       \{(\mM\vomega)^\trans\vx^+-\alpha\}
 \leq\alpha(\lambda-\lambda^+).
\]
Initialization and the cooling bank therefore have the same constants.
The nonnegative penalty bounds primal mass because
$(\mM\vomega)^\trans\vx\geq\alpha\vomega^\trans\vx$.

The restricted lazy walk is nonnegative and column substochastic, which
suffices for selected signed-flow propagation. Its moving comparator
remains monotone. For the analytical core use the \emph{restricted}
obstacle optimum at $r/2$. Its support volume is at most $2/r$, since
its nonnegative residual has weighted mass at most $\alpha$.
This core need not equal the full-graph optimum when $r/2<\rho$.
The zero-row slack comparison still gives $\lambda d_i/2$ off the core.
Consequently the entire weighted telescope and kinetic charge survive.

For the frozen energy, $\vomega^\trans\vt\leq1$ follows by summing
$\mM\vt=\vb$. Inverse positivity and the restricted residual bound give
$0\leq\vt-\vu_r\leq r\vomega$ and
$\|\vt-\vu_r\|_2^2\leq r$. The same freeze and rounding constants
follow. In the rounded energy estimates, the changed positive linear
penalty decreases under downward errors, and $|\mM|\vomega\leq\vomega$
supplies both weighted contractions. Neighbor sums in the implementation
include only neighbors in $A$ but their denominators and the safe rebase
subtraction still use $d_i$. The deficit proof therefore remains valid.
Finally the terminal optimum is the full $\vx_\rho^*$, so the original
full-graph RPPR bias proves the semantic output statement.
\end{proof}

The cutoff stays fixed for the whole run. Intermediate primal vectors are
not certified subsolutions, so residual-mass arguments for a dynamically
shrinking cutoff do not follow from this proposition.
